% 独立章节文件：可单独编辑，统一由 main.tex 引入。
\chapter{温补验证、研发结论与实施边界}
\section{常温和温补的验证应分开，但数据路径必须衔接}
\noindent\begin{tabularx}{\linewidth}{@{}p{34mm}YY@{}}\toprule
阶段 & 主要验证问题 & 必须具备的证据\\\midrule
常温算法验证 & 模型可辨识、独立工况误差、轴定义正确 & 24 ACC / 32 GYRO 工况、原始帧、反馈、报告\\
常温工艺验证 & 不同 PCB 与重装后重复性、装夹误差、转台不确定度 & GR\&R、跨板数据、计量证书与正式门限\\
温度模型研发 & 参数随温度变化、独立温度插值与补偿收益 & 至少拟合/验证分温度、热平衡、原始温度曲线\\
温度工艺验证 & 冷热往返滞回、上电漂移、重复循环、批次一致性 & 多次循环、冷/热两方向、返常温复测\\
整机验证 & 风机塔筒真实工况下倾角准确性、响应与振动鲁棒性 & 整机姿态融合、动态参考试验、温度和振动联合试验\\\bottomrule\end{tabularx}

\section{当前软件能够证明与不能证明的事情}
\textbf{能够计算：} 常温 A/B/C 的相同数据独立工况残差；温补在拟合温度下的独立工况残差；在未参与拟合温度点的插值残差；同一温补模型固定在 25°C 的对照残差。\par
\textbf{不能仅凭以上证明：} 全温区出厂合格、整机动态倾角达到特定竞品指标、全量程非线性、冷热滞回、多个 PCB 批次 GR\&R、2FDT 编码器差分速率的计量准确度及完整无人值守安全性。

\section{转为正式量产前的分级验证建议}
\begin{enumerate}
\item \textbf{台架接口确认（必须）：} 单独测试 PLC PV/SP/报警链路、2FDT 两轴方向、找零、停止和硬件急停；从低速空载开始，再按已批准的夹具方案加载。
\item \textbf{采样正确性确认（必须）：} 用有量值追踪的转台反馈和已知方向数据，确认 ACC/GYRO 轴符号及温度同步；特别复核曾见到的 GYRO 灵敏度系数数量级问题。
\item \textbf{重复性研究（建议）：} 同板多次常温标定、复装、不同操作者；再开展多板温补，分析参数分布和残差上界。
\item \textbf{冷热滞回（必须）：} 在升温和降温路径复测相同温度点，尤其返回 25°C 后重新执行关键 ACC/GYRO 独立验证工况。
\item \textbf{产品级验收（必须）：} 由目标动态倾角误差预算反推 ACC/GYRO 常温与温补门限，评审通过后再开启生产放行机制。
\end{enumerate}

\section{研发候选目标示例（\textbf{非当前软件出厂门限}）}
下面仅作为下一版研发讨论起点：ACC 独立验证方向 RMSE $\le0.02\sim0.03^\circ$；GYRO 低速矢量 RMSE $\le0.01\sim0.02^\circ/\mathrm s$；另研究低速静止漂移和 ACC/GYRO 相对轴系。它们\textbf{不是}如洋或 ST 官方验收阈值，也没有证据显示本系统跨板、跨温区已经稳定满足。正式阈值应由误差预算和实测不确定度推导。
